\begin{abstract}
Accurate prediction of horizontal-axis wind turbine (HAWT) wakes is
essential for wind farm layout optimisation, energy-yield estimation and
load assessment, yet it remains a major challenge for affordable
simulation tools. Large Eddy Simulation (LES) reproduces the wake physics
reliably but at a prohibitive computational cost, whereas unsteady
Reynolds-Averaged Navier-Stokes (URANS) closures systematically
mis-predict the wake recovery because of the assumptions embedded in the
eddy-viscosity hypothesis. This work proposes a CFD-ML methodology that
uses machine learning to \emph{define}, rather than post-process, the
turbulence closure of a $k$-$\omega$ SST RANS solver for HAWT wakes.
From a single high-fidelity LES of the NREL-5MW turbine operating in a
turbulent atmospheric boundary layer (ABL), an LES-equivalent
eddy-viscosity field $\nu_{t,eq}$ is extracted and used as the regression
target. A K-means clustering algorithm first segments the flow into
physically distinct regions, and a Multi-Layer Perceptron (MLP) is then
trained to predict the local correction $\Delta\nu_t = \nu_{t,eq} -
\nu_{t,\mathrm{RANS}}$. The resulting closure is embedded directly into a
modified OpenFOAM PIMPLE solver and evaluated in a fully predictive
setting. By design, the closure is calibrated on this single reference
case and then transferred, without retraining, to grid resolutions and
inflow velocities never seen during training: this single-case-to-many
transfer is the core premise of the methodology, not a limitation of it.
The trained network reproduces the LES-equivalent eddy viscosity with a
validation coefficient of determination $R^2 = 0.9952$, and the
clustering step is shown to improve the reconstruction in the
near-terrain region. Once embedded in the solver, the ML-based closure
consistently outperforms the standard RANS model in both the near- and
far-wake and remains close to the LES reference at the unseen inflow
velocities of $8$ and $15\,\mathrm{m/s}$. Most notably, on ultra-coarse
grids the model accurately captures the post-breakdown wake recovery at a
computational cost about $1.25$ times lower than the reference LES, while
also reproducing the rotor forces and power coefficients. These results
indicate that a turbulence closure trained once on an affordable amount
of high-fidelity data can deliver LES-quality wake predictions at a cost
compatible with practical wind farm applications.
\end{abstract}

\keywords{Wind turbine wake, URANS, Machine learning, Turbulence modelling, LES, CFD, Clustering, Neural Networks}
